\subsection{Entailment and the profile quotient}
\label{sec:fitting-entailment-quotient}

The preceding quotient construction raises a natural question: can the profile-saturation condition itself be generated by closing the admissible property domain under Fitting-style extensional entailment? The classical comparison suggests caution. In the mechanized Fitting analysis, the \(\delta\)-filter is a family of positive extensions: it contains the universal extension, excludes the empty extension, is upward closed under supersets, and is closed under intersections \cite{BenzmuellerFuenmayor2019}. This is closure of the positivity filter, not closure of the ambient property domain. The present \(A2_F^+\) interface already captures the corresponding monotonicity of positive support under extensional entailment.

I therefore distinguish that source-grounded filter behavior from the stronger domain principle
\[
\begin{split}
  EntCl^{adm}:\quad
  &Adm(X)\land
  \forall w\,Entail_F^+(w,X,Y)\\
  &\qquad\Longrightarrow Adm(Y).
\end{split}
\]
The stronger principle is unsuitable once FDE bottom is admitted. Since \(\bot_F\) has no positive members,
\[
  Entail_F^+(w,\bot_F,Y)
\]
holds vacuously for every world \(w\) and every rigid extension \(Y\). Lean consequently proves
\[
\boxed{
  EntCl^{adm}+Adm(\bot_F)
  \Longrightarrow
  \forall Y\,Adm(Y).
}
\]
Thus naive entailment closure of the admissible property universe reconstructs unrestricted bilateral comprehension. Combined with relevant exemplification consistency, the earlier comprehension obstruction returns immediately:
\[
\boxed{
  EntCl^{adm}+Adm(\bot_F)+CONS_{G,F}^{G,adm}
  \Longrightarrow
  \forall w,x\;\neg G_F^+(x,w).
}
\]
The distinction between filter closure and domain closure is therefore substantive rather than terminological.

A finite separation strengthens this negative result. Over two entities I take the complete \(4^2=16\) bilateral extension universe as admissible. Both entities actually exist at one reflexive world, and positive support is assigned exactly to extensions whose positive channel contains both entities. This model satisfies \(A2_F^+\), full domain-level entailment closure, closure under the basic FDE extension algebra, and quotient-compatible actual existence. Positivity nevertheless gives the two entities the same positive-property profile while an admissible singleton extension distinguishes them. Hence
\[
\boxed{
  A2_F^+
  +EntCl^{adm}
  +\text{FDE algebra closure}
  +\text{profile-compatible existence}
  \not\Longrightarrow Sat_P^{adm}
}
\]
in this explicit finite model. Entailment closure therefore does not generate the missing quotient structure.

The positive result takes a different form. For each world \(w\), define a bilateral profile closure on arbitrary rigid extensions by
\[
\begin{aligned}
  +Sat_w(Y)(x)
  &\Longleftrightarrow
  \exists y\,(x\approx_w^+y\land +Y(y)),\\
  -Sat_w(Y)(x)
  &\Longleftrightarrow
  \exists y\,(x\approx_w^+y\land -Y(y)).
\end{aligned}
\]
Lean proves that \(Sat_w\) is extensive and monotone and is idempotent up to bilateral extensional equivalence. Every \(Sat_w(Y)\) respects the positive-profile quotient. More strongly, if \(Z\) already respects the quotient and \(Y\) is support-included in \(Z\), then
\[
  Sat_w(Y)\preceq Z.
\]
Thus \(Sat_w(Y)\) is the least quotient-respecting bilateral extension above \(Y\) in the support-inclusion order.

This yields an exact fixed-point characterization of the earlier domain condition:
\[
\boxed{
  Sat_P^{adm}
  \Longleftrightarrow
  \forall w,Y\,
  \bigl(Adm(Y)\rightarrow Sat_w(Y)\equiv Y\bigr).
}
\]
The quotient route can therefore be read as selecting admissible properties from the fixed points of a canonical closure operator induced by positive-property observational equivalence.

Actualist quantification introduces one further boundary. Fitting entailment quantifies only over entities that actually exist at accessible worlds. Even if two entities lie in the same positive-profile class, the existence predicate can still distinguish them. I isolate the compatibility condition
\[
\begin{split}
  ExSat_P:\quad
  wRz\land x\approx_w^+y
  \Longrightarrow
  \bigl(E(z,x)\leftrightarrow E(z,y)\bigr).
\end{split}
\]
Under this premise Lean proves
\[
\boxed{
  Entail_F^+(w,X,Y)
  \Longrightarrow
  Entail_F^+(w,Sat_w(X),Sat_w(Y)).
}
\]
Hence actualist Fitting entailment descends through the profile closure when actual existence itself factors through the same quotient.

A second finite fixture shows why the existence premise matters. Two entities have the same positive-property profile, but at the sole accessible target world only one actually exists. An extension whose only positive member is the non-actual representative then entails an empty-positive extension vacuously. Profile closure transfers that positive membership to the equivalent actual representative, after which the saturated entailment fails. Thus the existence-compatibility condition is not merely a notational convenience, although the finite witness does not establish its global necessity for every alternative quotient construction.

The resulting architecture is more precise than the initial entailment-closure hypothesis. Positivity-filter closure remains governed by \(A2_F^+\); the admissible property universe should not be closed indiscriminately under entailment; profile saturation is instead a fixed-point condition for a canonical closure operator; and actualist entailment is compatible with that quotient only when actual existence respects the relevant profile classes. This relocates the next open problem from property-domain entailment closure to the principled status of the actualist quotient condition and to the algebraic/filter structure of the closure fixed points.
